\uuid{TGOG}
\exo7id{1824}
\titre{exo7 1824}
\auteur{drutu}
\organisation{exo7}
\datecreate{2003-10-01}
\isIndication{false}
\isCorrection{false}
\chapitre{Extremum, extremum lié}
\sousChapitre{Extremum, extremum lié}
\cpge{mpsi-mp2i.fonctions-deux-variables.b, pcsi.fonctions-deux-variables.b, ptsi.fonctions-deux-variables.b, bcpst1.fonctions-deux-variables.c, mp-mpi.fonctions-vectorielles.e}
\module{Calcul différentiel}
\niveau{L3}
\difficulte{}

\contenu{
\texte{
Ecriver la formule de Taylor de second ordre pour chacune des
fonctions suivantes au point $(x_0,y_0)$ donn\'e.
}
\begin{enumerate}
    \item \question{$f(x,y)=\sin (x+2y),\; (x_0,y_0)=(0,0)$ ;}
    \item \question{$f(x,y)=\frac{1}{x^2+y^2+1},\; (x_0,y_0)=(0,0)$ ;}
    \item \question{$f(x,y)=e^{-x^2-y^2}\cos xy,\; (x_0,y_0)=(0,0)$ ;}
    \item \question{$f(x,y)=\sin (xy)+\cos (xy),\; (x_0,y_0)=(0,0)$ ;}
    \item \question{$f(x,y)=e^{(x-1)^2}\cos y,\; (x_0,y_0)=(1,0)$.}
\end{enumerate}
}
